% fig09_testtime_compute.tex — 图9 Test-time Compute 收益曲线
% 《深度学习与大语言模型：前沿技术全景报告》图示库
% 由 dl_llm_report.tex 抽取，经 % fig09_testtime_compute.tex — 图9 Test-time Compute 收益曲线
% 《深度学习与大语言模型：前沿技术全景报告》图示库
% 由 dl_llm_report.tex 抽取，经 % fig09_testtime_compute.tex — 图9 Test-time Compute 收益曲线
% 《深度学习与大语言模型：前沿技术全景报告》图示库
% 由 dl_llm_report.tex 抽取，经 % fig09_testtime_compute.tex — 图9 Test-time Compute 收益曲线
% 《深度学习与大语言模型：前沿技术全景报告》图示库
% 由 dl_llm_report.tex 抽取，经 \input{figs/fig09_testtime_compute} 引用（caption/label 在主文件）
% 注意：本文件为片段，依赖主文件导言区的 tikz/pgfplots 宏包、颜色定义与 tikzset 样式，不可单独编译。

\begin{tikzpicture}
\begin{semilogxaxis}[width=0.8\textwidth, height=7.2cm, clip=false,
  xlabel={思考 token 预算（对数刻度）}, ylabel={任务准确率（示意）},
  xtick=\empty, ytick=\empty, axis lines=left, xmin=1, xmax=100000, ymin=0, ymax=100]
\addplot[very thick, cblue, domain=1:100000, samples=200] {8 + 80*(1 - exp(-(ln(x))^1.5/28))};
\addplot[thick, cgray, dashed, domain=1:100000] {90};
\node[lab, anchor=south west, align=left] at (axis cs:2,92.8) {任务难度上界（可望不可即）};
\node[lab, anchor=south east, align=left] at (axis cs:95000,82) {饱和区：成本翻倍、收益骤减};
\node[lab, anchor=west, align=left] at (axis cs:60,10) {近似幂律区间：\\预算翻倍 $\to$ 准确率可观提升};
\end{semilogxaxis}
\end{tikzpicture}
 引用（caption/label 在主文件）
% 注意：本文件为片段，依赖主文件导言区的 tikz/pgfplots 宏包、颜色定义与 tikzset 样式，不可单独编译。

\begin{tikzpicture}
\begin{semilogxaxis}[width=0.8\textwidth, height=7.2cm, clip=false,
  xlabel={思考 token 预算（对数刻度）}, ylabel={任务准确率（示意）},
  xtick=\empty, ytick=\empty, axis lines=left, xmin=1, xmax=100000, ymin=0, ymax=100]
\addplot[very thick, cblue, domain=1:100000, samples=200] {8 + 80*(1 - exp(-(ln(x))^1.5/28))};
\addplot[thick, cgray, dashed, domain=1:100000] {90};
\node[lab, anchor=south west, align=left] at (axis cs:2,92.8) {任务难度上界（可望不可即）};
\node[lab, anchor=south east, align=left] at (axis cs:95000,82) {饱和区：成本翻倍、收益骤减};
\node[lab, anchor=west, align=left] at (axis cs:60,10) {近似幂律区间：\\预算翻倍 $\to$ 准确率可观提升};
\end{semilogxaxis}
\end{tikzpicture}
 引用（caption/label 在主文件）
% 注意：本文件为片段，依赖主文件导言区的 tikz/pgfplots 宏包、颜色定义与 tikzset 样式，不可单独编译。

\begin{tikzpicture}
\begin{semilogxaxis}[width=0.8\textwidth, height=7.2cm, clip=false,
  xlabel={思考 token 预算（对数刻度）}, ylabel={任务准确率（示意）},
  xtick=\empty, ytick=\empty, axis lines=left, xmin=1, xmax=100000, ymin=0, ymax=100]
\addplot[very thick, cblue, domain=1:100000, samples=200] {8 + 80*(1 - exp(-(ln(x))^1.5/28))};
\addplot[thick, cgray, dashed, domain=1:100000] {90};
\node[lab, anchor=south west, align=left] at (axis cs:2,92.8) {任务难度上界（可望不可即）};
\node[lab, anchor=south east, align=left] at (axis cs:95000,82) {饱和区：成本翻倍、收益骤减};
\node[lab, anchor=west, align=left] at (axis cs:60,10) {近似幂律区间：\\预算翻倍 $\to$ 准确率可观提升};
\end{semilogxaxis}
\end{tikzpicture}
 引用（caption/label 在主文件）
% 注意：本文件为片段，依赖主文件导言区的 tikz/pgfplots 宏包、颜色定义与 tikzset 样式，不可单独编译。

\begin{tikzpicture}
\begin{semilogxaxis}[width=0.8\textwidth, height=7.2cm, clip=false,
  xlabel={思考 token 预算（对数刻度）}, ylabel={任务准确率（示意）},
  xtick=\empty, ytick=\empty, axis lines=left, xmin=1, xmax=100000, ymin=0, ymax=100]
\addplot[very thick, cblue, domain=1:100000, samples=200] {8 + 80*(1 - exp(-(ln(x))^1.5/28))};
\addplot[thick, cgray, dashed, domain=1:100000] {90};
\node[lab, anchor=south west, align=left] at (axis cs:2,92.8) {任务难度上界（可望不可即）};
\node[lab, anchor=south east, align=left] at (axis cs:95000,82) {饱和区：成本翻倍、收益骤减};
\node[lab, anchor=west, align=left] at (axis cs:60,10) {近似幂律区间：\\预算翻倍 $\to$ 准确率可观提升};
\end{semilogxaxis}
\end{tikzpicture}
